\intro{}


\section{Introduzione}

\defn{Ricorsione}{
    La ricorsione è un metodo di risoluzione di problemi computazionali in cui la soluzione di un problema dipende dalla soluzione di istanze più piccole dello stesso problema 
    e questo processo si realizza attraverso una funzione che richiama se stessa dall'interno del proprio corpo.
}

Una funzione ricorsiva è corretta e completa è sempre composta da esattamente due componenti fondamentali e inscindibili:

\begin{itemize}
    \item Caso base: il caso più elementare del problema, per cui la risposta è nota e immediata  la funzione lo risolve senza richiamare se stessa. È la condizione di arresto che garantisce la terminazione della catena ricorsiva. Senza di esso, la ricorsione continuerebbe all'infinito.
    \item Passo ricorsivo: il caso generale, in cui il problema non è ancora nella sua forma più semplice La funzione lo riduce in una versione più piccola e si richiama su quella versione ridotta, avanzando progressivamente verso il caso base.
\end{itemize}
\subsection{Struttura formale}
Matematicamente, una definizione ricorsiva si esprime sempre in questa forma:
\[
\begin{cases}
    \text{soluzione diretta} \ \text{se } n=\text{caso base}\\
    g\bigl(n,f(n')\bigr) \qquad \quad \text{altrimenti , dove } n'<n
\end{cases}
\]

Dove $n'$ è una versione ridotta di $n$, e $g$ è una qualsiasi operazione che combina il risultato della chiamata ricorsiva con il valore corrente. L'importante è che ogni passo ricorsivo garantisca che $n'$ converga verso il  caso base.


Da ora consideriamo la ricorsione non come “trucco”, ma come un insieme di pattern mentali ricorrenti.

\section{Ricorsione lineare}
\defn{Ricorsone lineare}{
    Una ricorsione è \textbf{lineare} quando ogni chiamataata genera esattamente una chiamata ricorsiva. La struttura è una catena: il problema di taglia n viene ridotto a un problema della stessa natura ma più piccolo, fino al caso base
}
Il principio è identico all'induzione matematica: il caso base corrisponde al passo base dell'induzione il passo ricorsivo corrisponde al passo induttivo

Non va simulata la catena nella testa.  Si Assume che la funzione funzioni già per input più piccoli.

\newpage

La ricorsione lineare si divide in tre sottotipi a seconda di dove avviene il lavoro rispetto alla chiamata ricorsiva.

\subsection{Post-order}
    Il lavoro avviene \textbf{dopo} la chiamata ricorsiva. La chiamata ricorsiva viene fatta prima, il suo risultato viene usato dopo.


\begin{codeexample}[Somma elementi positivi]
  \strategia{
    \begin{itemize}
      \item Facciamo la \textbf{ricorsione sull'indice} \inlinecode{i},
            avanzato di 1 ad ogni passo.
      \item \textbf{Ipotesi induttiva:} la funzione è corretta per qualsiasi
            indice strettamente maggiore di \inlinecode{i}.
      \item \textbf{Passo ricorsivo:} uso la funzione (corretta per ipotesi
            su \inlinecode{i+1}) e aggiungo il contributo di \inlinecode{v.at(i)}.
            \vspace{0.1cm}

            \inlinecode{int contributo = v.at(i) > 0 ? v.at(i) : 0;}
            \vspace{0.1cm}

            \inlinecode{return contributo + somma(v, i + 1);}

      \item \textbf{Caso base:} quando \inlinecode{i == (int)v.size()}
            non ci sono più elementi da esaminare; la somma è 0.
            \vspace{0.1cm}

            \inlinecode{if (i == (int)v.size()) return 0;}
    \end{itemize}
  }

  \begin{codebox}
    \begin{lstlisting}
int somma(const vector<int>& v, int i = 0) {
    if (i == (int)v.size()) return 0;
    int contributo = v.at(i) > 0 ? v.at(i) : 0;
    return contributo + somma(v, i + 1);
}
    \end{lstlisting}
  \end{codebox}

\end{codeexample}



\begin{codeexample}[Massimo]
  \strategia{
    \begin{itemize}
      \item Facciamo la \textbf{ricorsione sull'indice} \inlinecode{i},
            avanzato di 1 ad ogni passo.
      \item \textbf{Ipotesi induttiva:} la funzione è corretta per qualsiasi
            indice strettamente maggiore di \inlinecode{i}.
      \item \textbf{Passo ricorsivo:} uso la funzione (corretta per ipotesi
            su \inlinecode{i+1}) e confronto \inlinecode{v.at(i)} con il
            massimo del resto già noto.
            \vspace{0.1cm}

            \inlinecode{return max(v.at(i), massimo(v, i + 1));}

      \item \textbf{Caso base:} quando \inlinecode{i == (int)v.size() - 1}
            c'è un solo elemento, il massimo è lui.
            \vspace{0.1cm}

            \inlinecode{if (i == (int)v.size() - 1) return v.at(i);}
    \end{itemize}
  }

  \begin{codebox}
    \begin{lstlisting}
int massimo(const vector<int>& v, int i = 0) {
    if (i == (int)v.size() - 1) return v.at(i);
    return max(v.at(i), massimo(v, i + 1));
}
    \end{lstlisting}
  \end{codebox}

\end{codeexample}

  \subsection{Pre-order}
Il lavoro avviene prima della chiamata ricorsiva. L'effetto è l'opposto della post-order: si elabora in ordine naturale invece che inverso.
  \begin{codeexample}[Stampa da 1 a n]
  \strategia{
    \begin{itemize}
      \item Facciamo la \textbf{ricorsione sull'indice} \inlinecode{i},
            avanzato di 1 ad ogni passo.
      \item \textbf{Ipotesi induttiva:} la funzione è corretta per qualsiasi
            indice strettamente maggiore di \inlinecode{i}.
      \item \textbf{Passo ricorsivo:} uso la funzione (corretta per ipotesi
            su \inlinecode{i+1}) e stampo \inlinecode{i} prima di delegare
            il resto, il lavoro avviene alla discesa (pre-order).
            \vspace{0.1cm}

            \inlinecode{cout << i << " ";}
            \vspace{0.1cm}

            \inlinecode{stampa(i + 1, n);}

      \item \textbf{Caso base:} quando \inlinecode{i > n} non c'è più
            nulla da stampare.
            \vspace{0.1cm}

            \inlinecode{if (i > n) return;}
    \end{itemize}
  }

  \begin{codebox}
    \begin{lstlisting}
void stampa(int i, int n) {
    if (i > n) return;
    cout << i << " ";
    stampa(i + 1, n);
}
    \end{lstlisting}
  \end{codebox}

\end{codeexample}




















\begin{codeexample}[Somma con accumulatore]
  \strategia{
    \begin{itemize}
      \item Facciamo la \textbf{ricorsione sull'indice} \inlinecode{i},
            avanzato di 1 ad ogni passo, con \inlinecode{acc} che porta
            la somma parziale degli elementi positivi già visitati.
      \item \textbf{Ipotesi induttiva:} la funzione è corretta per qualsiasi
            indice strettamente maggiore di \inlinecode{i}.
      \item \textbf{Passo ricorsivo:} uso la funzione (corretta per ipotesi
            su \inlinecode{i+1}) aggiornando l'accumulatore con il contributo
            dell'elemento corrente  nessun calcolo rimane pendente al ritorno.
            \vspace{0.1cm}

            \inlinecode{int contributo = v.at(i) > 0 ? v.at(i) : 0;}
            \vspace{0.1cm}

            \inlinecode{return somma(v, i + 1, acc + contributo);}

      \item \textbf{Caso base:} quando \inlinecode{i == v.size()} tutti gli
            elementi sono stati processati; viene restituito direttamente
            l'accumulatore \inlinecode{acc}.
            \vspace{0.1cm}

            \inlinecode{if (i == v.size()) return acc;}
    \end{itemize}
  }

  \begin{codebox}
    \begin{lstlisting}
int somma(const vector<int>& v, int i = 0, int acc = 0) {
    if (i == v.size()) return acc;
    int contributo = v.at(i) > 0 ? v.at(i) : 0;
    return somma(v, i + 1, acc + contributo);
}
    \end{lstlisting}
  \end{codebox}

\end{codeexample}
  \subsection{Tail recursion}
  La chiamata ricorsiva è l'ultima operazione assoluta della funzione nessun calcolo pendente dopo di essa. Lo stato che nella post-order si accumula nello stack viene qui portato in avanti come parametro accumulatore.

  \begin{codeexample}[Somma con accumulatore]
  \strategia{
    \begin{itemize}
      \item Facciamo la \textbf{ricorsione sull'indice} \inlinecode{i},
            avanzato di 1 ad ogni passo, con \inlinecode{acc} che porta
            la somma parziale degli elementi positivi già visitati.
      \item \textbf{Ipotesi induttiva:} la funzione è corretta per qualsiasi
            indice strettamente maggiore di \inlinecode{i}.
      \item \textbf{Passo ricorsivo:} uso la funzione (corretta per ipotesi
            su \inlinecode{i+1}) aggiornando l'accumulatore con il contributo
            dell'elemento corrente  nessun calcolo rimane pendente al ritorno.
            \vspace{0.1cm}

            \inlinecode{int contributo = v.at(i) > 0 ? v.at(i) : 0;}
            \vspace{0.1cm}

            \inlinecode{return somma(v, i + 1, acc + contributo);}

      \item \textbf{Caso base:} quando \inlinecode{i == (int)v.size()} tutti
            gli elementi sono stati processati; viene restituito direttamente
            l'accumulatore \inlinecode{acc}.
            \vspace{0.1cm}

            \inlinecode{if (i == (int)v.size()) return acc;}
    \end{itemize}
  }

  \begin{codebox}
    \begin{lstlisting}
int somma(const vector<int>& v, int i = 0, int acc = 0) {
    if (i == (int)v.size()) return acc;
    int contributo = v.at(i) > 0 ? v.at(i) : 0;
    return somma(v, i + 1, acc + contributo);
}
    \end{lstlisting}
  \end{codebox}

\end{codeexample}

\begin{codeexample}[Massimo comun divisore di Euclide]
  La proprietà matematica su cui si fonda l'algoritmo è la seguente:
  \[
    \text{MCD}(a,\ b) = \text{MCD}(b,\ a \bmod b)
  \]
  \textbf{Perché è vera.} Dalla divisione euclidea si ha $a = q \cdot b + r$
  con $r = a \bmod b$. Qualsiasi divisore comune $d$ di $a$ e $b$ divide anche
  $r = a - q \cdot b$, e viceversa. I due insiemi di divisori comuni di
  $(a,b)$ e $(b,r)$ sono quindi identici, e in particolare lo è il loro massimo.

  \strategia{
    \begin{itemize}
      \item \textbf{Ricorsione su:} $b$, ridotto via $a \bmod b$ ad ogni passo.
            Poiché $0 \le a \bmod b < b$, il secondo argomento decresce
            strettamente la sequenza raggiunge $b = 0$ in tempo finito.
      \item \textbf{Ipotesi induttiva:} la funzione è corretta per qualsiasi
            valore di $b$ strettamente minore di quello corrente.
      \item \textbf{Passo ricorsivo:} per la proprietà sopra,
            $\text{MCD}(a,b) = \text{MCD}(b,\ a \bmod b)$;
            per ipotesi la funzione è già corretta su $(b,\ a \bmod b)$,
            quindi la chiamata ricorsiva restituisce il risultato corretto.
            La chiamata è \textit{tail}: $a \bmod b$ viene calcolato prima
            di essere passato, nessun calcolo rimane pendente al ritorno.
      \item \textbf{Caso base:} quando $b = 0$, ogni intero divide $0$,
            quindi $\text{MCD}(a,\ 0) = a$.
    \end{itemize}
  }

  \begin{codebox}
    \begin{lstlisting}
int mcd(int a, int b) {
    if (b == 0) return a;    
    return mcd(b, a % b);   
}
    \end{lstlisting}
  \end{codebox}

\end{codeexample}



Ora andiamo a vedere il secondo pattern la \textbf{Tail recursion}.

\section{Tail recursion}
\defn{Tail recursion}{
    Una chiamata ricorsiva è in coda, quando è l'ultima operazione assoluta della funzione nessun calcolo rimane pendente dopo di essa. Il valore restituito della chiama ricorsiva viene restituito direttamente ,senza ulteriori elaborazioni.
}

La tail recursion si divide in due famiglie principali: con accumulatore e senza accumulatore:

\subsection{Senza accumulatore}
La chiamata ricorsiva è  l'ultima operazione, ma non c'è nulla da comunicare in avanti ,  il risultato della chiamata ricorsiva viene restituito direttamente senza modifiche.

Si usa quando la ricorsione \textbf{elimina} possibilità invece di costruire qualcosa. Tipicamente ricerche, verifiche, riduzioni immediate. 


\begin{codeexample}[Ricerca lineare]
  \strategia{
    \begin{itemize}
      \item Facciamo la \textbf{ricorsione sull'indice} \inlinecode{i},
            avanzato di 1 ad ogni passo.
      \item \textbf{Ipotesi induttiva:} la funzione è corretta per qualsiasi
            indice strettamente maggiore di \inlinecode{i}.
      \item \textbf{Passo ricorsivo:} se \inlinecode{v.at(i)} non è il
            target, uso la funzione (corretta per ipotesi su \inlinecode{i+1})
            e ne restituisco direttamente il risultato  nessun calcolo
            rimane pendente al ritorno.
            \vspace{0.1cm}

            \inlinecode{return contiene(v, target, i + 1);}

      \item \textbf{Casi base:} due condizioni distinte.
            \vspace{0.1cm}

            \inlinecode{if (i == (int)v.size()) return false;}
            \vspace{0.1cm}

            \inlinecode{if (v.at(i) == target) return true;}
    \end{itemize}
  }

  \begin{codebox}
    \begin{lstlisting}
bool contiene(const vector<int>& v, int target, int i = 0) {
    if (i == (int)v.size()) return false;
    if (v.at(i) == target)  return true;
    return contiene(v, target, i + 1);
}
    \end{lstlisting}
  \end{codebox}

\end{codeexample}









\begin{codeexample}[Vettore ordinato]
  \strategia{
    \begin{itemize}
      \item Facciamo la \textbf{ricorsione sull'indice} \inlinecode{i},
            avanzato di 1 ad ogni passo.
      \item \textbf{Ipotesi induttiva:} la funzione è corretta per qualsiasi
            indice strettamente maggiore di \inlinecode{i}.
      \item \textbf{Passo ricorsivo:} se \inlinecode{v.at(i) <= v.at(i+1)},
            il confronto corrente è soddisfatto; uso la funzione (corretta
            per ipotesi su \inlinecode{i+1}) e ne restituisco direttamente
            il risultato  nessun calcolo rimane pendente al ritorno.
            \vspace{0.1cm}

            \inlinecode{return ordinato(v, i + 1);}

      \item \textbf{Casi base:} due condizioni distinte.
            \vspace{0.1cm}

            \inlinecode{if (v.at(i) > v.at(i + 1)) return false;}
            \vspace{0.1cm}

            \inlinecode{if (i == (int)v.size() - 1) return true;}
    \end{itemize}
  }

  \begin{codebox}
    \begin{lstlisting}
bool ordinato(const vector<int>& v, int i = 0) {
    if (i == (int)v.size() - 1)  return true; 
    if (v.at(i) > v.at(i + 1))   return false;
    return ordinato(v, i + 1);
}
    \end{lstlisting}
  \end{codebox}

\end{codeexample}


\subsection{Con accumulatore}
Serve quando si deve costruire o calcolare qualcosa lungo il percorso.
Questo sottotipo a sua volta ha altri sottotipi:

\subsubsection{Accumulatore singolo numerico}
Lo stato è un singolo valore numerico che si aggiorna ad ogni passo.




\begin{codeexample}[Fattoriale con accumulatore]

  Il \textbf{fattoriale} di un intero non-negativo $n$ ammette una
  definizione ricorsiva diretta:
  \[
    n! = \begin{cases}
      1              & \text{se } n = 0 \\
      n \cdot (n-1)! & \text{se } n > 0
    \end{cases}
  \]
  Per ottenere una \textbf{tail recursion}, si introduce un accumulatore
  \inlinecode{acc} che porta il prodotto parziale in avanti invece di
  lasciarlo pendente sullo stack. Poiché \inlinecode{acc} è un dettaglio
  implementativo interno, si separa la logica in due funzioni: una
  \textbf{funzione helper} che espone \inlinecode{acc} come parametro
  esplicito, e una \textbf{funzione wrapper} pubblica che la chiama
  inizializzando \inlinecode{acc} al valore neutro della moltiplicazione,
  ovvero \inlinecode{1}.

  \strategia{
    \begin{itemize}
      \item Facciamo la \textbf{ricorsione su} \inlinecode{n},
            ridotto di 1 ad ogni passo, con accumulatore \inlinecode{acc}
            che porta il prodotto parziale.
      \item \textbf{Ipotesi induttiva:} la funzione è corretta per qualsiasi
            valore di \inlinecode{n} strettamente minore di quello corrente.
      \item \textbf{Passo ricorsivo:} aggiorno \inlinecode{acc}
            moltiplicando per \inlinecode{n} e delego  nessun calcolo
            rimane pendente al ritorno.
            \vspace{0.1cm}

            \inlinecode{return fact\_helper(n - 1, acc * n);}

      \item \textbf{Caso base:} quando \inlinecode{n == 0} il calcolo
            è completo; viene restituito direttamente \inlinecode{acc}.
            \vspace{0.1cm}

            \inlinecode{if (n == 0) return acc;}
    \end{itemize}
  }

  \begin{codebox}
    \begin{lstlisting}
int fact_helper(int n, long long acc) {
    if (n == 0) return acc;
    return fact_helper(n - 1, acc * n);
}

int fact(int n) {
    return fact_helper(n, 1);   // acc inizializzato al valore neutro
}
    \end{lstlisting}
  \end{codebox}

\end{codeexample}

\subsection{ Accumulatore stringa}
Lo stato è una stringa che si costruisce progressivamente.


\begin{codeexample}[Inversione stringa con accumulatore]

  Per ottenere una tail recursion, si introduce un accumulatore
  \inlinecode{acc} di tipo \inlinecode{string} che costruisce la
  stringa invertita progressivamente, portandola in avanti invece
  di lasciarla pendente sullo stack. \inlinecode{acc} viene
  inizializzato alla stringa vuota, valore neutro della
  concatenazione.

  \strategia{
    \begin{itemize}
      \item Facciamo la \textbf{ricorsione sull'indice} \inlinecode{i},
            avanzato di 1 ad ogni passo, con accumulatore \inlinecode{acc}
            che porta la stringa invertita parziale.
      \item \textbf{Ipotesi induttiva:} la funzione è corretta per qualsiasi
            indice strettamente maggiore di \inlinecode{i}.
      \item \textbf{Passo ricorsivo:} prependo \inlinecode{s.at(i)}
            all'accumulatore e delego  nessun calcolo rimane
            pendente al ritorno.
            \vspace{0.1cm}

            \inlinecode{return inverti(s, i + 1, s.at(i) + acc);}

      \item \textbf{Caso base:} quando \inlinecode{i == (int)s.size()}
            tutti i caratteri sono stati processati; viene restituito
            direttamente \inlinecode{acc}.
            \vspace{0.1cm}

            \inlinecode{if (i == s.size()) return acc;}
    \end{itemize}
  }

  \begin{codebox}
    \begin{lstlisting}
string inverti(const string& s, int i = 0, string acc = "") {
    if (i == (int)s.size()) return acc;
    return inverti(s, i + 1, s.at(i) + acc);
}
    \end{lstlisting}
  \end{codebox}

\end{codeexample}
\subsection{ Accumulatori multipli paralleli}
Serve quando lo stato richiede più valori che evolvono insieme. I parametri accumulatori sono due o più e si aggiornano in parallelo ad ogni passo.

















\begin{codeexample}[Fibonacci con accumulatori paralleli]

  La successione di Fibonacci è definita come:
  \[
    F(n) = \begin{cases}
      0               & \text{se } n = 0 \\
      1               & \text{se } n = 1 \\
      F(n-1) + F(n-2) & \text{se } n > 1
    \end{cases}
  \]
  La definizione standard genera una ricorsione multipla con
  sottoproblemi sovrapposti. Per ottenere una tail recursion si
  introducono due accumulatori paralleli \inlinecode{a} e
  \inlinecode{b} che portano avanti gli ultimi due valori della
  sequenza, aggiornandosi ad ogni passo senza calcoli pendenti
  al ritorno.

  \strategia{
    \begin{itemize}
      \item Facciamo la \textbf{ricorsione su} \inlinecode{n},
            ridotto di 1 ad ogni passo, con due accumulatori
            paralleli \inlinecode{a} e \inlinecode{b} che portano
            rispettivamente $F(k)$ e $F(k+1)$ per il passo corrente.
      \item \textbf{Ipotesi induttiva:} la funzione è corretta per
            qualsiasi valore di \inlinecode{n} strettamente minore
            di quello corrente.
      \item \textbf{Passo ricorsivo:} aggiorno entrambi gli
            accumulatori in parallelo  \inlinecode{a} diventa
            \inlinecode{b} e \inlinecode{b} diventa \inlinecode{a+b}
             e delego. Nessun calcolo rimane pendente al ritorno.
            \vspace{0.1cm}

            \inlinecode{return fib(n - 1, b, a + b);}

      \item \textbf{Caso base:} quando \inlinecode{n == 0},
            \inlinecode{a} contiene già $F(0)$; viene restituito
            direttamente.
            \vspace{0.1cm}

            \inlinecode{if (n == 0) return a;}
    \end{itemize}
  }

  \begin{codebox}
    \begin{lstlisting}
int fib_helper(int n, long long a, long long b) {
    if (n == 0) return a;
    return fib_helper(n - 1, b, a + b);
}

    int fib(int n) {
    return fib_helper(n, 0, 1);   // a=F(0)=0, b=F(1)=1
}
    \end{lstlisting}
  \end{codebox}

\end{codeexample}













\section{Divide et impera}
\defn{}{
Divide et impera è una strategia ricorsiva che risolve un problema dividendolo in sotto problemi indipendenti della stessa natura, risolvendo ognuno per ipotesi induttiva,e poi combinando i risultati parziali in una soluzione completa.
}

Nel pattern Divide et impera ci sono due sottotipi, i quali sono:
\subsection{Divide et impera con combine esplicita}
Il lavoro principale avviene nella fase combina, dopo che entrambe la metà sono state risolte, si fa un lavoro non banale per unire i  risultati. 


\begin{codeexample}[Merge sort]








  Il \textbf{Merge Sort} è un algoritmo di ordinamento \textit{divide et impera}.
  Opera in due fasi:
  \begin{enumerate}
    \item \textbf{Divide:} trova il punto medio $m = \lfloor(l+r)/2\rfloor$ e
          richiama ricorsivamente \inlinecode{mergeSort} sulle due metà
          \inlinecode{[l,m]} e \inlinecode{[m+1,r]}.
    \item \textbf{Merge:} fonde le due metà ordinate confrontando elemento
          per elemento e copiando ogni volta il minore in un array di supporto,
          poi ricopia il risultato in \inlinecode{arr[l..r]}.
  \end{enumerate}


\begin{center}

\begin{tikzpicture}[
    nodes={minimum width=0.7cm, minimum height=0.7cm},
    row sep=-\pgflinewidth,
    column sep=-\pgflinewidth,
    >=stealth
]

% Livello 0: Array originale
\matrix(orig)[matrix of nodes,
    row 1/.style={nodes={draw=none}},
    nodes={draw}]
{
    \tiny{0} & \tiny{1} & \tiny{2} & \tiny{3}\\
    5 & 2 & 4 & 1\\
};

% Livello 1: Split
\matrix(left)[matrix of nodes,
    row 1/.style={nodes={draw=none}},
    nodes={draw},
    below left=1.2cm and 0.5cm of orig]
{
    \tiny{0} & \tiny{1}\\
    5 & 2\\
};

\matrix(right)[matrix of nodes,
    row 1/.style={nodes={draw=none}},
    nodes={draw},
    below right=1.2cm and 0.5cm of orig]
{
    \tiny{2} & \tiny{3}\\
    4 & 1\\
};

% Livello 2: Merge (ordinato)
\matrix(mleft)[matrix of nodes,
    row 1/.style={nodes={draw=none}},
    nodes={draw, fill=green!15},
    below left=1.2cm and 0.5cm of $(left)!0.5!(right)$]
{
    \tiny{0} & \tiny{1}\\
    2 & 5\\
};

\matrix(mright)[matrix of nodes,
    row 1/.style={nodes={draw=none}},
    nodes={draw, fill=green!15},
    below right=1.2cm and 0.5cm of $(left)!0.5!(right)$]
{
    \tiny{2} & \tiny{3}\\
    1 & 4\\
};

% Livello 3: Array finale ordinato
\matrix(final)[matrix of nodes,
    row 1/.style={nodes={draw=none}},
    nodes={draw, fill=orange!20},
    below=1.2cm of $(mleft)!0.5!(mright)$]
{
    \tiny{0} & \tiny{1} & \tiny{2} & \tiny{3}\\
    1 & 2 & 4 & 5\\
};

% Frecce split
\draw[->, thick, red!70!black] (orig-2-2.south) -- (left-1-1.north)
    node[midway, left, font=\scriptsize] {split};
\draw[->, thick, red!70!black] (orig-2-3.south) -- (right-1-2.north);

% Frecce merge
\draw[->, thick, green!50!black] (left-2-2.south) -- (mleft-1-1.north)
    node[midway, left, font=\scriptsize] {merge};
\draw[->, thick, green!50!black] (right-2-1.south) -- (mright-1-2.north);

% Frecce merge finale
\draw[->, thick, orange!70!black] (mleft-2-2.south) -- (final-1-2.north)
    node[midway, left, font=\scriptsize] {merge};
\draw[->, thick, orange!70!black] (mright-2-1.south) -- (final-1-3.north);

\end{tikzpicture}


\end{center}






    \strategia{
        \begin{itemize}
        \item Facciamo la \textbf{ricorsione sugli estremi} \inlinecode{l} e \inlinecode{r} dimezzando l'intervallo ad oggni passo. 
        \item \textbf{Ipotesi induttiva:}: la funzione è corretta per qualsiasi intervallo di dimensione strettamente minore di \inlinecode{r-l+1}
        \item \textbf{Passo ricorsivo:} per ipotesi le due metà sono già ordinate.
        \item \textbf{Caso base:} quando \inlinecode{l>=r} c'è un solo elemento ed è già ordinato
        \end{itemize}
    }



    \begin{codebox}
    \begin{lstlisting}
    void merge(vector<int>& v, int l, int m, int r) {
    vector<int> sx, dx;

    for (int i = l; i <= m; i++) sx.push_back(v.at(i));
    for (int i = m + 1; i <= r; i++) dx.push_back(v.at(i));

    int i = 0, j = 0, k = l;
    while (i < (int)sx.size() && j < (int)dx.size())
        v.at(k++) = (sx.at(i) <= dx.at(j)) ? sx.at(i++) : dx.at(j++);
    while (i < (int)sx.size()) v.at(k++) = sx.at(i++);
    while (j < (int)dx.size()) v.at(k++) = dx.at(j++);
}
    \end{lstlisting}
    \end{codebox}






    \begin{codebox}
    \begin{lstlisting}
    void mergeSort(vector<int>& v, int l, int r) {
    if (l >= r) return;
    int m = l + (r - l) / 2;     // DIVIDE
    mergeSort(v, l, m);           // CONQUER sinistra
    mergeSort(v, m + 1, r);       // CONQUER destra
    merge(v, l, m, r);            // COMBINE
}
    \end{lstlisting}
    \end{codebox}
\end{codeexample}


\begin{codeexample}[Massimo]
    \strategia{
    \begin{itemize}
    \item Facciamo la \textbf{riscorsione sugli estremi} \inlinecode{l} e \inlinecode{r}, dimmezando l'intervallo.
    \item \textbf{Ipotesi induttiva:} la funzione è coretta èer qualsiasi intervallo di dimensione strettamente minore di \inlinecode{r-l+1}.
    \item \textbf{Passo ricorsivo:} per ipotesi ho il massimo delle due metà; prendo il maggiore tra i due. 
    \item \textbf{Caso base:} quando \inlinecode{l==r} cìè un solo elemento 
    \end{itemize}
    }
    \begin{codebox}
    \begin{lstlisting}
    int massimo(const vector<int>& v, int l, int r){
    if(l==r) return v.at(l);
    int m = l+(r-1)/2;
    int maxSx = massimo(v, l, m);        // CONQUER sinistra
    int maxDx = massimo(v, m + 1, r);    // CONQUER destra
    return max(maxSx, maxDx);           // Combina 
    }
    \end{lstlisting}
    \end{codebox}
\end{codeexample}




\subsection{Divide et impera senza combina, con riduzione diretta}
Non c'è fase di combine  dopo aver diviso, si procede su una sola delle parti. È un caso degenere di Divide et impera che assomiglia alla ricorsione lineare, ma il criterio di divisione è $n/2$ invece di $ n-1$.

\begin{codeexample}[Fast power]
\strategia{
    \begin{itemize}
        \item Facciamo la \textbf{ricorsione} su \inlinecode{exp}, dimezzato ad ogni passo.
        \item \textbf{Ipotesi induttiva:} la funzione è coretta per qualsiasi valore di \inlinecode{exp} strettamente minore di quello corrente 
        \item \textbf{Passo ricorsivo:} per ipotesi ho \inlinecode{base\^(exp/2)}.
        \item \textbf{Caso base:}\ilinecode{exp== 0}
    \end{itemize}
}

    \begin{codebox}
    \begin{lstlisting}
    int fastPow(int base,int exp){
     if(exp== 0) return 1; 
     int half = fastPow(base,exp/2);
      if(exp%2 == 0) return half *hallf;
      else
            return base*half*half;
    }
    \end{lstlisting}
    \end{codebox}
\end{codeexample}

\begin{codeexample}
    Suppponiamo che qualcuno pensi a un numero tra 1 e 100, e noi dobbiamo indovinarlo. Ci verrà detto solo se è troppo alto , troppo basso, o se il numero è esatto.
    Per fare ciò ci sono due strategie:
    \begin{itemize}
    \item \textbf{Strategia  non ottimale}: Proviamo 1, poi 2, poi 3$\ldots$ se il numero fosse 99, ci vorebbero 99 tentativi. Ogni tentativo elimina un solo numero.
    \item \textbf{Strategia ottimale:} Ad ogni tentativo puntiamo sempre al centro dell'intervallo:
    \begin{enumerate}
        \item Proviamo $50 \rightarrow$ troppo alto elimino i numeri da 50 a 100.
        \item Proviamo $25 \rightarrow$ troppo alto elimino i numeri da 25 a 49.
        \item Proviamo $12 \rightarrow$ troppo basso limino i numeri da 1 a 12 .
        \item Proviamo $18 \rightarrow$ troppo basso limino i numeri da 13 a 18.
        \item Proviamo $21 \rightarrow$ troppo basso limino i numeri da 19 a 21.
        \item 23 trovato.
    \end{enumerate}
    \end{itemize}
    


    \begin{center}
        \begin{tikzpicture}[
    nodes={minimum width=1cm, minimum height=0.7cm},
    row sep=-\pgflinewidth,
    column sep=-\pgflinewidth,
    >=stealth
]

% === Passo 0: Array 1..100, mid=50 ===
\matrix(s0)[matrix of nodes, nodes={draw}]
{
    1 & \dots & |[fill=yellow!40]| 50 & \dots & 100\\
};
\node[right=0.5cm of s0, font=\scriptsize, align=left]
    {mid = 50\\ 23 < 50 $\Rightarrow$ vai a sx};

% === Passo 1: Array 1..49, mid=25 ===
\matrix(s1)[matrix of nodes, nodes={draw},
    below=1.2cm of s0]
{
    1 & \dots & |[fill=yellow!40]| 25 & \dots & |[fill=red!15]| 49\\
};
\node[right=0.5cm of s1, font=\scriptsize, align=left]
    {mid = 25\\ 23 < 25 $\Rightarrow$ vai a sx};

% === Passo 2: Array 1..24, mid=12 ===
\matrix(s2)[matrix of nodes, nodes={draw},
    below=1.2cm of s1]
{
    1 & \dots & |[fill=yellow!40]| 12 & \dots & |[fill=red!15]| 24\\
};
\node[right=0.5cm of s2, font=\scriptsize, align=left]
    {mid = 12\\ 23 > 12 $\Rightarrow$ vai a dx};

% === Passo 3: Array 13..24, mid=18 ===
\matrix(s3)[matrix of nodes, nodes={draw},
    below=1.2cm of s2]
{
    |[fill=red!15]| 13 & \dots & |[fill=yellow!40]| 18 & \dots & 24\\
};
\node[right=0.5cm of s3, font=\scriptsize, align=left]
    {mid = 18\\ 23 > 18 $\Rightarrow$ vai a dx};

% === Passo 4: Array 19..24, mid=21 ===
\matrix(s4)[matrix of nodes, nodes={draw},
    below=1.2cm of s3]
{
    |[fill=red!15]| 19 & \dots & |[fill=yellow!40]| 21 & \dots & 24\\
};
\node[right=0.5cm of s4, font=\scriptsize, align=left]
    {mid = 21\\ 23 > 21 $\Rightarrow$ vai a dx};

% === Passo 5: Array 22..24, mid=23 → TROVATO ===
\matrix(s5)[matrix of nodes, nodes={draw},
    below=1.2cm of s4]
{
    22 & |[fill=green!40]| 23 & 24\\
};
\node[right=0.5cm of s5, font=\scriptsize, align=left, text=green!50!black]
    {mid = 23\\ \textbf{Trovato!}};

% === Frecce tra i passi ===
\draw[->, thick] (s0.south) -- (s1.north);
\draw[->, thick] (s1.south) -- (s2.north);
\draw[->, thick] (s2.south) -- (s3.north);
\draw[->, thick] (s3.south) -- (s4.north);
\draw[->, thick] (s4.south) -- (s5.north);

% === Titolo e legenda ===
\node[above=0.4cm of s0, font=\bfseries]
    {Ricerca Binaria --- Cerco il 23 in [1, 100]};

\node[below=0.6cm of s5, font=\scriptsize, align=center]
    {
    \colorbox{yellow!40}{~mid~} = elemento centrale \quad
    \colorbox{red!15}{~scartato~} \quad
    \colorbox{green!40}{~trovato!~}\\[3pt]
    };

\end{tikzpicture}
    \end{center}
    \begin{codebox}
    \begin{lstlisting}
    bool ricercaBinaria(const vector<int>& v, int target, int da, int a){
        if(da>=a){
        int m = (da+a)/2
        if(v.at(m==2)
            return true;
        else 
            if(v.at(m)>e)
                return ricercaBinaria(v,e,da,m-1);
            else
                return ricercaBinaria(v,e,m+1,a);
        }
        else
            return false;
    }
    \end{lstlisting}
    \end{codebox}
\end{codeexample}

\section{Backtracking}
\defn{Backtracking}{
Il backtracking è una tecnica ricorsiva per esplorare sistematicamente uno spazio di scelte costruendo una soluzione incrementalmente. Ad ogni passo si effettua una scelta, si assume per ipotesi induttiva che la ricorsione esplori correttamente tutto ciò che segue, poi si annulla la scelta per esplorare le alternative.
}
Nel backtracking ci sono tre sottotipi i quali sono: 

\subsection{Backtracking su scelte binarie}
Ad ogni passo la scelta è binaria:includo l'elemento corrente oppure no. 

\begin{codeexample}[sottoinsiemi]
    \strategia{
        \begin{itemize}
        \item Facciamo la \textbf{ricorsione sull'iindice } \inlinecode{i}, avanzando ad ogni passo.
        \item \textbf{Passo ricorsivo:} er ogni elemento \inlinecode{v.at(i)} esploro i due rami — incluso e non incluso — usando la funzione corretta per ipotesi su \inlinecode{i+1}.
        \item \textbf{Caso base:} \inlinecode{i == v.size()} la soluzione pariziale è completa, e viene registrata.


        \end{itemize}
        }
    \begin{lstlisting}
    void subsets(const vector<int> & v, int i, vector<int> & corrente,vector<vector<int>> & res){
    if(i== v.size()){
        res.push_back(corrente);
        return;
    }
    subsets(v,i+1,corrente,res);
    corrente.push_back(v.at(i));
     subsets(v, i + 1, corrente, res);
      corrente.pop_back();   
    }
    \end{lstlisting}
    
\end{codeexample}

\subsection{ Backtracking su permutazioni}
Ad ogni passo le scelte sono filtrate da vincoli che dipendono dallo stato corrente.


\begin{codeexample}[Permutazioni]

  \strategia{
    \begin{itemize}
      \item Facciamo la \textbf{ricorsione sulla lunghezza} di
            \inlinecode{corrente}, aumentata di 1 ad ogni passo.
      \item \textbf{Ipotesi induttiva:} la funzione è corretta per
            qualsiasi stato con un elemento in più già scelto in
            \inlinecode{corrente}.
      \item \textbf{Passo ricorsivo:} per ogni elemento non ancora
            usato, lo aggiungo a \inlinecode{corrente}, uso la funzione
            (corretta per ipotesi) per esplorare il resto, poi annullo
            la scelta per provare le alternative.
            \vspace{0.1cm}

            \inlinecode{usato.at(i) = true;}
            \vspace{0.1cm}

            \inlinecode{corrente.push\_back(v.at(i));}
            \vspace{0.1cm}

            \inlinecode{permutazioni(v, usato, corrente, res);}
            \vspace{0.1cm}

            \inlinecode{corrente.pop\_back();}
            \vspace{0.1cm}

            \inlinecode{usato.at(i) = false;}

      \item \textbf{Caso base:} quando
            \inlinecode{(int)corrente.size() == (int)v.size()}
            la permutazione è completa; viene registrata.
            \vspace{0.1cm}

            \inlinecode{if ((int)corrente.size() == (int)v.size())}
            \inlinecode{res.push\_back(corrente);}
    \end{itemize}
  }

  \begin{codebox}
    \begin{lstlisting}
void permutazioni(const vector<int>& v, vector<bool>& usato,
                  vector<int>& corrente, vector<vector<int>>& res) {
    if ((int)corrente.size() == (int)v.size()) {
        res.push_back(corrente);
        return;
    }
    for (int i = 0; i < (int)v.size(); i++) {
        if (!usato.at(i)) {
            usato.at(i) = true;
            corrente.push_back(v.at(i));            
            permutazioni(v, usato, corrente, res);  
            corrente.pop_back();                    
            usato.at(i) = false;                  
        }
    }
}
    \end{lstlisting}
  \end{codebox}

\end{codeexample}







\subsection{ Backtracking con vincoli strutturali }
Ad ogni passo le scelte sono filtrate da vincoli che dipendono dallo stato corrente. 



\begin{codeexample}[Parentesi ben formate]

  \strategia{
    \begin{itemize}
      \item Facciamo la \textbf{ricorsione sulla lunghezza} di
            \inlinecode{corrente}, aumentata di 1 ad ogni passo,
            con contatori \inlinecode{open} e \inlinecode{close}
            che tracciano le parentesi già inserite.
      \item \textbf{Ipotesi induttiva:} la funzione è corretta per
            qualsiasi stato con un carattere in più già aggiunto a
            \inlinecode{corrente}.
      \item \textbf{Passo ricorsivo:} provo ad aggiungere
            \inlinecode{(} se \inlinecode{open < n}, e \inlinecode{)}
            se \inlinecode{close < open}; il pruning è implicito nelle
            condizioni. L'undo è implicito poiché \inlinecode{corrente}
            è passata per valore.
            \vspace{0.1cm}

            \inlinecode{if (open < n)}
            \vspace{0.1cm}

            \inlinecode{\ \ \ \ parentesi(n, open+1, close, corrente+"(", res);}
            \vspace{0.1cm}

            \inlinecode{if (close < open)}
            \vspace{0.1cm}

            \inlinecode{\ \ \ \ parentesi(n, open, close+1, corrente+")", res);}

      \item \textbf{Caso base:} quando
            \inlinecode{(int)corrente.size() == 2*n}
            la stringa è completa; viene registrata.
            \vspace{0.1cm}

            \inlinecode{if ((int)corrente.size() == 2*n)}
            \inlinecode{res.push\_back(corrente);}
    \end{itemize}
  }

  \begin{codebox}
    \begin{lstlisting}
void parentesi(int n, int open, int close,
               string corrente, vector<string>& res) {
    if ((int)corrente.size() == 2 * n) {
        res.push_back(corrente);
        return;
    }
    if (open < n)
        parentesi(n, open + 1, close, corrente + "(", res);
    if (close < open)
        parentesi(n, open, close + 1, corrente + ")", res);
}
    \end{lstlisting}
  \end{codebox}

\end{codeexample}